\documentclass{article}

\usepackage[utf8]{inputenc}
\usepackage[italian]{babel}
\usepackage{amsmath}
\usepackage{amssymb}
\usepackage{amsthm}
\usepackage{booktabs}
\usepackage{hyperref}

% Importa il pacchetto tensors dal livello superiore
\usepackage{../tensors}

% Ridefinizione di \L come nel documento principale (Thesis.tex)
\renewcommand{\L}{\mathbb{L}}

\title{Corrispondenza tra Ordinamento dei Tensori e Layout di Memoria delle Matrici}
\author{}
\date{}

\begin{document}

\maketitle

Questo documento descrive formalmente la corrispondenza tra le nozioni tensoriali di \emph{ordinamento naturale} e \emph{ordinamento inverso} (definiti nel capitolo 1 della tesi) e i layout di memoria standard per matrici bidimensionali, noti come \emph{column-major} e \emph{row-major order}.

\section{Formalizzazione per Matrici (Tensori 2-dimensionali)}

Sia data una matrice $\Mx{A} \in \mathbb{R}^{n_1 \times n_2}$, indicizzata dalla coppia di indici $(i_1, i_2)$, dove $i_1 \in [n_1]$ indica la riga e $i_2 \in [n_2]$ indica la colonna. 

\subsection{Ordinamento Naturale e Column-Major}
L'ordinamento naturale $\L \colon \mdom! \to [\prod_{j=1}^d n_j]$ definisce la vettorizzazione scorrendo per prima la prima dimensione. Per $d=2$:
\begin{itemize}
    \item I passi (o strides) sono definiti come $s_1 = 1$ e $s_2 = n_1$.
    \item La mappa di linearizzazione è:
    \[
    \L(i_1, i_2) = i_1 + n_1 (i_2 - 1).
    \]
\end{itemize}
Poiché il passo associato all'indice di riga $i_1$ è $s_1 = 1$, tale indice varia più velocemente. In memoria, gli elementi adiacenti appartengono alla stessa colonna:
\[
\L(1,1) = 1, \quad \L(2,1) = 2, \quad \dots, \quad \L(n_1,1) = n_1, \quad \L(1,2) = n_1 + 1, \quad \dots
\]
Questo corrisponde esattamente al layout \textbf{column-major order} (tipico di Fortran, MATLAB, R, Julia).
\begin{itemize}
    \item \textbf{Vettorizzazione:} L'operatore $\Vec(\Mx{A})$ basato su $\L$ effettua un incolonnamento per colonne (ad esempio, \texttt{A(:)} in MATLAB o \texttt{A.flatten('F')} in Python).
\end{itemize}

\subsection{Ordinamento Inverso e Row-Major}
L'ordinamento inverso (reverse ordering) $\L^* \colon \mdom! \to [\prod_{j=1}^d n_j]$ inverte l'ordine dei passi. Per $d=2$:
\begin{itemize}
    \item I passi sono definiti come $s_2^* = 1$ e $s_1^* = n_2$.
    \item La mappa di linearizzazione è:
    \[
    \L^*(i_1, i_2) = i_2 + n_2 (i_1 - 1).
    \]
\end{itemize}
Poiché il passo associato all'indice di colonna $i_2$ è $s_2^* = 1$, tale indice varia più velocemente. In memoria, gli elementi adiacenti appartengono alla stessa riga:
\[
\L^*(1,1) = 1, \quad \L^*(1,2) = 2, \quad \dots, \quad \L^*(1,n_2) = n_2, \quad \L^*(2,1) = n_2 + 1, \dots
\]
Questo corrisponde esattamente al layout \textbf{row-major order} (tipico di C, C++, Python/NumPy/PyTorch per impostazione predefinita).
\begin{itemize}
    \item \textbf{Vettorizzazione:} L'operazione di vettorizzazione associata a $\L^*$ corrisponde a scorrere la matrice riga per riga (ad esempio, \texttt{A.flatten('C')} in Python).
\end{itemize}

\section{Estensione a Tensori multidimensionali ($d > 2$)}

Sia $\Tn{X} \in \Rmsiz!$ un tensore d-dimensionale indicizzato da $(\miwc) \in \mdom!$.

\begin{description}
    \item[Ordinamento Naturale ($\L$):] I passi sono $s_1 = 1$, $s_{j+1} = s_j n_j = n_1 \cdots n_j$. L'indice più a sinistra ($i_1$) varia più velocemente in memoria, mentre l'indice più a destra ($i_d$) varia più lentamente. Corrisponde a un layout multidimensionale di tipo \textbf{column-major} (Fortran).
    \item[Ordinamento Inverso ($\L^*$):] I passi sono $s_d^* = 1$, $s_{j-1}^* = s_j^* n_j = n_j \cdots n_d$. L'indice più a destra ($i_d$) varia più velocemente in memoria, mentre l'indice più a sinistra ($i_1$) varia più lentamente. Corrisponde a un layout multidimensionale di tipo \textbf{row-major} (C, Python).
\end{description}

\section{Tabella di Confronto}

\begin{table}[ht]
\centering
\begin{tabular}{@{}lllll@{}}
\toprule
Ordinamento & Strides ($d=2$) & Layout Matrice & Indice Veloce & Standard Informatici \\ \midrule
Naturale ($\L$) & $s_1=1, s_2=n_1$ & Column-major & $i_1$ (riga) & MATLAB, Julia, R, Fortran \\
Inverso ($\L^*$) & $s_2^*=1, s_1^*=n_2$ & Row-major & $i_2$ (colonna) & C, C++, Python (NumPy) \\ \bottomrule
\end{tabular}
\caption{Confronto tra ordinamenti tensoriali e layout di memoria matriciali.}
\end{table}

\end{document}
